Bei einem adiabatischen Prozess wird keine Wärme ausgetauscht, $Q = 0$. Die vom Gas verrichtete Arbeit geht vollständig auf Kosten der inneren Energie:
\begin{align}
 W_\text{Gas} &= -\Delta U = \frac{f}{2}\,nR\left(T_1-T_2\right)
\end{align}
Wasserstoff (H$_2$) ist zweiatomig, $f = \fHO$. Die Temperaturdifferenz ist $T_1-T_2 = \dT$:
\begin{align}
 W_\text{Gas} &= \WGF \\
   &= \frac{\fHO}{2}\cdot\nWs\cdot\ncR\cdot\dT \\
   &= \WG \approx \result{\WGP{3}{2}}
\end{align}
